% Conclusion: no new results. Each answer cites the table that carries it.
\section{Summary of findings}
\paragraph{RQ1: is exact CT supervision enough?} No. On rendered LIDC-IDRI data with exact truth, lesion insertion lowered absorption a little
at no cost in bone error, and the pre-specified hypothesis H1 passed (Table~\ref{tab:v1_h1}). On real radiographs, however, the v1 models
absorbed more than twice what the reference model absorbed, and on the locked VinDr-RibCXR audit the v1 model still absorbs 0.289 of an added
lesion. Exact supervision of the bone image does not by itself keep lesions out of it.

\paragraph{RQ2: does the twin regulariser help without losing fidelity?} Yes, for bone-region fidelity. Absorption falls with the twin weight
(H1$'$, Table~\ref{tab:dose_response}); the final ensemble absorbs 0.051 of an added lesion on VinDr-RibCXR against 0.149 for the released
reference model and 0.178 for the matched control (H3$'$, Table~\ref{tab:audit_contrasts}), while its bone-region PSNR is non-inferior to both
(H2$'$, Table~\ref{tab:h2}). Structural similarity to the dual-energy image missed its margin, mainly through fine texture, and on exact-truth
data the regulariser raises the bone error. The ablation shows that the effect needs a consistency term confined to the lesion and gradients
through both twins (Table~\ref{tab:ablation_twin}), and it appears with every encoder tried (Table~\ref{tab:encoders}).

\paragraph{RQ3: does the audit measure something real?} In part. Against the matched control, real-nodule contrast rose consistently with the
regulariser on both nodule data sets, which agrees with the audit. Against the reference model the answer depends on the data set: higher
contrast on JSRT, lower on Node21, where the reference model suppresses more strongly (Table~\ref{tab:nodule_cnr}). The audit and the evidence
without insertion agree on the effect of the regulariser; they do not agree on a single ranking of different models.

\paragraph{RQ4: does it matter downstream, and is the output reliable?} The soft-tissue image used alone did not improve classification or
detection, and adding it to the original image gave small gains (Tables~\ref{tab:downstream_cls} and~\ref{tab:downstream_det}). The
ensemble spread ranks the bone error well on rendered data, and conformal intervals reach their nominal coverage there, but neither holds on
real radiographs (Table~\ref{tab:reliability}); an exploratory recalibration on 24 real pairs restored the coverage of the intervals
(Table~\ref{tab:bone_recal}). On paediatric radiographs the suppressor was non-inferior for pneumonia classification, kept its absorption
ordering, and its uncertainty rose 2.7-fold (Chapter~\ref{ch:external}).

\paragraph{RQ5: do classifiers gain from the suppressor's images?} Not overall. A multi-stream classifier that read the radiograph,
soft-tissue and bone images did not improve macro-AUROC on any of three data sets, and it did not beat a control of the same capacity or
channel stacking (H6 and H7, Table~\ref{tab:cls_primary}). The largest effect on both adult data sets was a small gain for nodules and
masses, significant on NIH ChestX-ray14 only and offset there by a loss for diffuse findings (H8, Table~\ref{tab:cls_groups}); attribution
maps of the soft-tissue stream located nodules and masses more often than those of the radiograph stream (H10,
Table~\ref{tab:cls_pointing}). The suppressor's own uncertainty did not help a classifier decide when to abstain
(Table~\ref{tab:cls_selective}).

\paragraph{RQ6: does CT teach a Swin X-ray classifier?} Not with this design. Distilling a 2.5D Swin CT classifier into the X-ray model
left the classification of real radiographs unchanged: on the five shared findings the distilled classifier differed from one started from
ImageNet by $-0.0001$ $[-0.0016, 0.0014]$ in mean AUROC (H11), and neither pretraining on the projections nor the teacher on top of it made
a difference (H12, Table~\ref{tab:xm_hyp}); on the projections and in zero-shot transfer the distilled models were slightly worse
(Table~\ref{tab:xm_transfer}). The CT and X-ray heat maps of the same patient did not agree (H14, Table~\ref{tab:xm_crossmodal}). The Swin
classifier was slightly less accurate than the DenseNet-121 of Chapter~\ref{ch:classifier} (H13), almost entirely because of hernia; per
class it lost clearly on hernia and slightly on infiltration and matched it on the other 12 diseases (Table~\ref{tab:xm_per_class}). It was far less sensitive to image noise than DenseNet-121, and under rib emphasis the distilled classifier
flipped slightly fewer finding-free images than the one started from ImageNet (H15, $p=0.044$ uncorrected), as did the classifier pretrained
on the projections without the teacher (Table~\ref{tab:xm_halluc}). Grad-CAM++ of the Swin classifiers located expert boxes less than half
as often as that of DenseNet-121, while Integrated Gradients and decoder attention did about as well (Table~\ref{tab:xm_pointing}).

\section{Contributions}
The thesis contributes a counterfactual lesion audit that makes lesion loss by bone suppressors measurable on real radiographs, a twin
lesion-invariance regulariser that lowers it, a recipe that trains such suppressors on public data only, two pre-specified evaluations
on locked data with downstream tasks, an external validation and a full record of deviations and negative results, and a pre-specified test,
against a control of the same capacity, of whether classifiers gain from bone-suppressed inputs. It also contributes an explainable,
hallucination-aware Swin framework for chest CT and X-ray, with CT-to-X-ray distillation, a test of heat-map agreement between the
modalities and counterfactual hallucination tests, evaluated under a fourth pre-specified plan; its negative results bound what such a
transfer can add with public data of this size.

\section{Recommendations}
Whoever evaluates a bone suppressor should report how much of an inserted lesion it absorbs, on real radiographs and at sites shared by all
compared models, next to a measure of how much bone it removes, and not rely on PSNR and SSIM alone. Whoever uses one should show the
suppressed image beside the original rather than instead of it, should treat the ensemble spread as a warning that the image is unusual rather
than as a calibrated error bar, and should not use a suppressor on a population or image type it was not evaluated on. Whoever builds a
classifier should not expect bone-suppressed inputs to raise its overall accuracy, and should test any benefit per finding type against a
control of the same capacity. Whoever explains a Swin classifier of radiographs should prefer Integrated Gradients or decoder attention to
Grad-CAM++ and should score the maps against expert boxes; whoever claims that a classifier is robust should test it with perturbations that
add no disease, because accuracy did not reveal the nearly threefold difference in noise sensitivity between two encoders of almost equal
accuracy.

\section{Outlook}
The audit applies to any image-to-image model that is meant to remove one structure and keep another. The next steps are a reader study, other
lesion types, higher resolution, real dual-energy data for calibration, and paediatric training data. For the cross-modality framework, a
stable and stronger CT teacher comes first: the 2.5D recipe collapsed in one of seven seeds and was weak for lung lesions, so a full 3D
encoder and many more volumes are needed before distillation can be judged on its merits. Heat-map agreement should then be measured with
Integrated Gradients or attention maps rather than Grad-CAM++, and disagreement between the CT and X-ray models could be tested as a signal
for when to abstain.
